This dissertation presents the design, implementation and evaluation of a multiclient file server upon the seL4 Microkit framework, motivated by the aim of mitigating IPC overhead, a key performance constraint in microkernel-based systems, through asynchronous client-server communication. The system supports a rich set of FS operations, including namespace, data and metadata management, and multiclient access control.

The file server is structured around modular and flexible core abstractions including blocks, inodes, directories and file descriptors, providing clear separation between persistent FS metadata and per-client runtime state. Two communication models are implemented over this core system: a synchronous model based on protected procedure calls, and an asynchronous model based on notifications, shared memory queues and batched request processing. A transparent API abstracts communication details while preserving consistent operation semantics across both models.

The asynchronous model demonstrated reduced IPC overhead through batching and reduced synchronisation frequency, improving throughput and lowering average operation latency under suitable workloads. However, these benefits are workload-dependent, with limited gains under highly interdependent workloads. A key insight is that communication structure alone can significantly influence system behaviour independently of core FS design. The results further suggest that, although evaluated in a single-threaded environment, the asynchronous design could enable substantially greater concurrent processing on multicore systems through non-blocking execution and overlap of request generation, servicing and response handling.

The work demonstrates that complex, stateful services can be realised entirely in user space within the Microkit framework for seL4, and provides insight into trade-offs between synchronous and asynchronous communication for future microkernel-based system services.


